% Options for packages loaded elsewhere
\PassOptionsToPackage{unicode}{hyperref}
\PassOptionsToPackage{hyphens}{url}
\documentclass[
  ignorenonframetext,
]{beamer}
\newif\ifbibliography
\usepackage{pgfpages}
\setbeamertemplate{caption}[numbered]
\setbeamertemplate{caption label separator}{: }
\setbeamercolor{caption name}{fg=normal text.fg}
\beamertemplatenavigationsymbolsempty
% remove section numbering
\setbeamertemplate{part page}{
  \centering
  \begin{beamercolorbox}[sep=16pt,center]{part title}
    \usebeamerfont{part title}\insertpart\par
  \end{beamercolorbox}
}
\setbeamertemplate{section page}{
  \centering
  \begin{beamercolorbox}[sep=12pt,center]{section title}
    \usebeamerfont{section title}\insertsection\par
  \end{beamercolorbox}
}
\setbeamertemplate{subsection page}{
  \centering
  \begin{beamercolorbox}[sep=8pt,center]{subsection title}
    \usebeamerfont{subsection title}\insertsubsection\par
  \end{beamercolorbox}
}
% Prevent slide breaks in the middle of a paragraph
\widowpenalties 1 10000
\raggedbottom
\AtBeginPart{
  \frame{\partpage}
}
\AtBeginSection{
  \ifbibliography
  \else
    \frame{\sectionpage}
  \fi
}
\AtBeginSubsection{
  \frame{\subsectionpage}
}
\usepackage{iftex}
\ifPDFTeX
  \usepackage[T1]{fontenc}
  \usepackage[utf8]{inputenc}
  \usepackage{textcomp} % provide euro and other symbols
\else % if luatex or xetex
  \usepackage{unicode-math} % this also loads fontspec
  \defaultfontfeatures{Scale=MatchLowercase}
  \defaultfontfeatures[\rmfamily]{Ligatures=TeX,Scale=1}
\fi
\usepackage{lmodern}
\ifPDFTeX\else
  % xetex/luatex font selection
\fi
% Use upquote if available, for straight quotes in verbatim environments
\IfFileExists{upquote.sty}{\usepackage{upquote}}{}
\IfFileExists{microtype.sty}{% use microtype if available
  \usepackage[]{microtype}
  \UseMicrotypeSet[protrusion]{basicmath} % disable protrusion for tt fonts
}{}
\makeatletter
\@ifundefined{KOMAClassName}{% if non-KOMA class
  \IfFileExists{parskip.sty}{%
    \usepackage{parskip}
  }{% else
    \setlength{\parindent}{0pt}
    \setlength{\parskip}{6pt plus 2pt minus 1pt}}
}{% if KOMA class
  \KOMAoptions{parskip=half}}
\makeatother
\setlength{\emergencystretch}{3em} % prevent overfull lines
\providecommand{\tightlist}{%
  \setlength{\itemsep}{0pt}\setlength{\parskip}{0pt}}
\usepackage{bookmark}
\IfFileExists{xurl.sty}{\usepackage{xurl}}{} % add URL line breaks if available
\urlstyle{same}
\hypersetup{
  hidelinks,
  pdfcreator={LaTeX via pandoc}}

\author{\texorpdfstring{}{}}
\date{}

\begin{document}

\section{Diagrams as Generative
Models}\label{sec:diagrammatic-cognition}

\begin{frame}{Why the Brain Prefers Diagrams}
\protect\phantomsection\label{why-the-brain-prefers-diagrams}
We can now substantially strengthen our initial architectural claim from
\autoref{sec:introduction}: formal commutative diagrams do not merely
provide a convenient pedagogical representation illustrating abstract
case structure---they mathematically constitute the exact native
\emph{computational format} through which biological cognitive agents
actively maintain and query their internal generative models
representing relational environmental structure.

This strong structural claim draws support from powerful converging
empirical evidence:

\begin{enumerate}
\item
  \textbf{Computational advantage} (Larkin \& Simon,
  {[}-@larkin1987diagram{]}): Diagrams enable search, recognition, and
  inference operations that are computationally prohibitive in
  sentential format. A commutative case diagram allows the agent to
  verify consistency (does the direct path equal the composed path?) by
  simple spatial inspection.
\item
  \textbf{Free ride inferences} (Shimojima,
  {[}-@shimojima1996reasoning{]}): Properties of the diagram that are
  perceptually available but would require explicit computation in a
  sentential format. In a case diagram, transitivity of grammatical
  relations is \emph{visible}---the existence of a path from NOM to DAT
  through ACC is spatially apparent.
\item
  \textbf{Hybrid reasoning} (Giardino,
  {[}-@giardino2017diagrammatic{]}): Mathematical diagrams engage a mode
  of reasoning that combines perceptual pattern recognition with
  background theoretical knowledge. Case diagrams similarly engage both
  the perceptual system (spatial layout) and linguistic knowledge (case
  constraints, verb valency).
\item
  \textbf{Peirce's existential graphs}: Peirce's graphical logic system
  demonstrated that first-order logic can be conducted entirely
  diagrammatically, without algebraic symbols. Our case diagrams extend
  this tradition: the relational structure of a sentence is represented
  graphically, and inference proceeds by diagram manipulation
  (adding/removing nodes, composing morphisms).
\end{enumerate}
\end{frame}

\begin{frame}{P600 Signals and Garden-Path Reanalysis in Diagrammatic
Models}
\protect\phantomsection\label{p600-signals-and-garden-path-reanalysis-in-diagrammatic-models}
The standard predictive processing framework---for which active
inference operates as the most formally mathematically developed
version---provides a remarkably natural, mechanistically precise account
detailing exactly how biological agents actively deploy these
diagrammatic representations cognitively during real-time processing:

\begin{enumerate}
\item
  \textbf{Top-down structural predictions}: The currently active
  internal case diagram continuously generates precise, high-precision
  predictions anticipating incoming expected sensory input (e.g., the
  system computationally predicts ``a nominative-marked noun phrase must
  immediately appear because the parsed transitive verb structurally
  demands an active agent'').
\item
  \textbf{Bottom-up prediction errors}: Incoming sensory words that
  physically violate the model's top-down diagrammatic predictions
  instantly generate massive, measurable prediction errors (e.g.,
  encountering a structurally unexpected morphological case marker
  directly triggers a quantifiable P600 event-related neural potential
  measurable in the biological brain).
\item
  \textbf{Belief updating}: The diagram is updated to accommodate the
  prediction error, potentially restructuring the assignment of entities
  to case roles (garden-path reanalysis)
\item
  \textbf{Precision weighting}: The enriched weights on morphisms serve
  as \emph{precision parameters} that control the relative influence of
  prior expectations and incoming evidence. A high-weight morphism
  generates strong predictions that are costly to override; a low-weight
  morphism generates weak predictions that are easily overridden.
\end{enumerate}
\end{frame}

\begin{frame}{Three Falsifiable ERP Predictions}
\protect\phantomsection\label{three-falsifiable-erp-predictions}
The predictive processing account generates quantitative, falsifiable
predictions about neural responses to case-marking violations. In the
active inference framework, a case-assignment violation triggers a
\emph{prediction error} whose amplitude scales with the precision of the
violated expectation---which is precisely the enriched hom-value of the
violated morphism:

\begin{equation}
\text{PE}(f) \propto \pi_f \cdot |\mu_{\text{predicted}} - \mu_{\text{observed}}|
\label{eq:pe-precision-error}
\end{equation}

where \(\pi_f = \mathcal{C}(A,B)\) is the enriched weight (precision) of
the morphism \(f: A \to B\) and \(\mu\) are the expected vs.~observed
case features. This yields three concrete electrophysiological
predictions:

\begin{enumerate}
\item
  \textbf{P600 amplitude scales with morphism weight}: A case violation
  on a high-weight morphism (NOM→ACC in a transitive clause,
  \(w = 0.9\)) should elicit a larger P600 than a violation on a
  low-weight morphism (NOM→INS in an experiencer construction,
  \(w = 0.4\)). The ratio of P600 amplitudes should approximate the
  ratio of enriched weights.
\item
  \textbf{N400 reflects distributional expectation}: Semantic case
  violations---where the case-marked NP satisfies the morphological case
  but not the distributional proto-role requirements (e.g., an inanimate
  NOM in an agentive construction)---should elicit N400 effects
  proportional to the \emph{magnitude deficit} of the violated enriched
  category (\autoref{sec:enriched-categories}).
\item
  \textbf{Garden-path reanalysis costs track magnitude}: The processing
  cost of reanalyzing a garden-path sentence's case structure should
  correlate with the \emph{change in categorical magnitude} between the
  initial and revised case diagrams, since magnitude quantifies how much
  relational information the agent's generative model encodes.
\end{enumerate}
\end{frame}

\begin{frame}{Integration: Five Pillars Become One Generative Model}
\protect\phantomsection\label{integration-five-pillars-become-one-generative-model}
The full picture emerges when we combine all five pillars within the
active inference framework:

\begin{enumerate}
\tightlist
\item
  \textbf{Case categories} (\autoref{sec:case-systems}) provide the
  \emph{objects and morphisms} of the generative model---the vocabulary
  of roles and relations
\item
  \textbf{Categorial grammar} (\autoref{sec:categorial-grammar})
  provides the \emph{composition rules}---how roles combine to form
  structured derivations
\item
  \textbf{DisCoCat / DisCoCirc} (\autoref{sec:categorical-semantics},
  \autoref{sec:compact-closure-complexity},
  \autoref{sec:discocirc-discourse}) provides the \emph{semantic
  functor} and discourse extension---mapping syntactic structure to
  distributional meaning
\item
  \textbf{Enriched structure} (\autoref{sec:enriched-categories})
  provides the \emph{precision parameters}---graded weights that control
  inference
\item
  \textbf{Topos-theoretic bridges} (\autoref{sec:topos-theory}) provide
  \emph{transfer theorems}---ensuring consistency across formalizations
\end{enumerate}

The active inference agent uses this combined structure as a single,
integrated generative model. Each scenario it encounters---a sentence
heard, a scene observed, an action planned---is interpreted by
instantiating a case diagram from structure (1), parsing the input using
rules (2), computing meaning via the semantic functor (3), weighting
confidence using enriched structure (4), and transferring results across
representational formats using bridge techniques (5).

This is \emph{total cognitive scenario understanding}: the agent doesn't
just parse a sentence or assign case labels---it constructs a complete,
internally consistent, generic, strongly typed, dynamically updating
model of the relational structure of the situation, and uses that model
to predict, explain, and act.
\end{frame}

\begin{frame}[fragile]{804 Automated Tests Confirm the Formalism Is
Executable}
\protect\phantomsection\label{automated-tests-confirm-the-formalism-is-executable}
The framework developed in this paper is not merely theoretical---it is
computationally verified through a comprehensive implementation and test
suite that exercises every categorical construction discussed above.

\begin{block}{System Architecture and Categorical Core}
\protect\phantomsection\label{system-architecture-and-categorical-core}
The categorical core (\texttt{CaseCategory}, \texttt{EnrichedCategory},
\texttt{AlignmentFunctor}, \texttt{NaturalTransformation}) is
implemented in Python with set-based object tracking and list-based
morphism storage, enforcing categorical axioms at construction time. Six
additional modules extend the core: \texttt{FluidSFunctor}
(context-dependent alignment parameterized by volition),
\texttt{CaseDiagramBelief} and active inference computations
(variational free energy, prediction error, belief updating), the
\textbf{\texttt{src/daif/} subpackage} (7 modules, 24 symbols---full
distributional RL inference: push-forward returns, quantile TD, VMP,
Bethe FE, EIG, ERP profiles, policy selection, and metrics),
\texttt{CasePOVM} and quantum case assignment (POVM-based probability
via \autoref{eq:eq-8-1} in \autoref{sec:quantum-semantics}),
\texttt{DitransitiveSentence} (three-argument verb support), and
\texttt{CaseFrameValidator} (cognitive security via type-violation
detection). The visualization layer produces all manuscript figures
programmatically, ensuring exact correspondence between formal claims
and visual evidence. The DisCoPy integration library (version 1.2.2)
provides an independent validation path: pregroup types (\texttt{Ty}),
lexical entries (\texttt{Word}), cup contractions (\texttt{Cup}), cap
expansions (\texttt{Cap}), type permutations (\texttt{Swap}), normal
form computation (\texttt{normal\_form()}), and circuit depth analysis
(\texttt{depth()}) are exercised against the same categorical structures
described in \autoref{sec:categorial-grammar} and
\autoref{sec:categorical-semantics}.
\end{block}

\begin{block}{Automated Test Suite and Verification}
\protect\phantomsection\label{automated-test-suite-and-verification}
The implementation is validated by \textbf{804 automated tests} across
47 test files with \textbf{≥90\% code coverage} (enforced in the build
configuration). Every test uses real mathematical computations---no
mocks or fakes. The \textbf{per-category inventory} (counts, modules,
and DAIF file breakdown) is listed in
\autoref{sec:test-suite-inventory}.

This computational verification demonstrates that the category-theoretic
framework is not just a mathematical convenience but a \emph{working
computational architecture}---the categorical abstractions compile,
execute, and produce verifiable results, bridging the gap between formal
theory and implemented system.
\end{block}
\end{frame}

\end{document}
